\chapter{Whittaker 系数的保守性}\label{ch:13}
\srcpage{64}
本章旨在解释 Whittaker 系数为何能够检测自守范畴的温和部分\cite{FR25,Gai18}。在范畴 Langlands 的 de Rham 形式中，预期有
\[\DMod(\Bun_G)\simeq\IndCoh_{\mathrm{Nilp}}(\LocSys_{\checkG}),\]

\begin{transnote}{含糊点拨}
$\IndCoh_{\mathrm{Nilp}}$ 是 $\LocSys_{\check G}$ 上奇异支撑落在全局幂零锥内的 ind-凝聚层：幂零锥由 $\check{\mathfrak g}^*$ 的幂零锥沿 $\LocSys$ 拉回。de Rham 侧用 $\IndCoh$ 而不是 $\QCoh$，是因为 Hecke 特征层的谱侧天然带奇异支撑；只有温和部分支撑为零，才退化为 $\QCoh$。所以「温和 $\leftrightarrow$ $\QCoh$」是一条独立陈述，不是定义。[来源：nLab "ind-coherent sheaves"（该条目无 Last revised，2026-11-17 访问）；并见原书 p.64 本段与 [Gai18]、[FR25]。]
\end{transnote}
以及
\[\DMod(\Bun_G)^{\mathrm{temp}}\simeq\QCoh(\LocSys_{\checkG}).\]
在 Betti 理论或可构造理论中，相应的谱叠是 $\LocSys^{\mathrm{restr}}_{\checkG}$，而自守侧使用 $\Shv_{\mathrm{Nilp}}(\Bun_G)$，即层范畴中具有幂零奇异支撑的部分。本章记录的定理是重整化 Whittaker 系数函子的保守性：
\begin{theorem*}{}
\[\begin{aligned}
\DMod(\Bun_G)^{\mathrm{temp}}&\xrightarrow{\operatorname{Coeff}^{\mathrm{ren}}}\QCoh(\LocSys_{\checkG}),\\
\Shv_{\mathrm{Nilp}}(\Bun_G)^{\mathrm{temp}}&\xrightarrow{\operatorname{Coeff}^{\mathrm{ren}}_{\mathrm{restr}}}\QCoh(\LocSys^{\mathrm{restr}}_{\checkG}).
\end{aligned}\]
这里，“保守”意味着：若系数为零，则对象本身为零。
\end{theorem*}
在 de Rham 实现中，对应的可构造范畴是特征簇包含于全局幂零锥的正则完整 $D$ 模的 ind-完备化。它并非整个 $\DMod(\Bun_G)$，而是正则完整且具有幂零支撑的部分。

\section{普通系数函子与重整化系数函子}\label{sec:13.1}\index{Whittaker 系数!重整化|hyperpage}\index{余作用|hyperpage}
固定 Borel 子群 $B=TN$，以及非退化加性特征标 $\psi:N\to\Ga$。令 $\Bun_N^\Omega$ 为通常的 $\Omega$-扭转 $N$-丛叠，它带有投影 $p:\Bun_N^\Omega\to\Bun_G$ 和 Whittaker 特征标 $\psi:\Bun_N^\Omega\to\mathbb A^1$。忽略标准的上同调归一化，第一个 Whittaker 系数是紧支撑 de Rham 上同调
\begin{definition*}{}
\[\operatorname{coeff}(\cF):=C_{\mathrm{dR},c}(\Bun_N^\Omega,p^!(\cF)\otimes\psi^!(\exp))\in\operatorname{Vect}.\]
\end{definition*}
等价地，它是将 $\cF$ 拉回到 Whittaker 叠后，以指数层扭转，再 $!$-推前到一点。移位的选择与 Verdier 对偶和下文的微局部指标公式相容。

谱作用给出
\[\QCoh(\LocSys_{\checkG})\otimes\DMod(\Bun_G)\longrightarrow\DMod(\Bun_G).\]
通过对偶，一个标量值泛函 $c:\DMod(\Bun_G)\to\operatorname{Vect}$ 产生一个谱值泛函
\[\begin{split}
\DMod(\Bun_G)&\xrightarrow{\operatorname{coact}}\QCoh(\LocSys_{\checkG})\otimes\DMod(\Bun_G)\\
&\xrightarrow{\operatorname{Id}\otimes c}\QCoh(\LocSys_{\checkG}).
\end{split}\]
取 $c=\operatorname{coeff}$，便得到 $\operatorname{Coeff}^{\mathrm{ren}}$。它的余单位，也就是取全局截面后得到的标量形式，恢复原来的标量系数：
\[\Gamma(\LocSys_{\checkG},\operatorname{Coeff}^{\mathrm{ren}}(\cF))\simeq\operatorname{coeff}(\cF).\]
\begin{note}
单凭这一恢复性质，不能刻画取值于 $\QCoh(\LocSys_{\checkG})$ 的对象，因为 $\LocSys_{\checkG}$ 一般不是仿射的；其定义是上面的余作用公式。同样的说明也适用于受限 Betti／可构造情形。在谱范畴满足通常的刚性假设时，给出到 $\QCoh(\LocSys_{\checkG})$ 的 $\QCoh(\LocSys_{\checkG})$-线性函子，等价于给出取值于 $\operatorname{Vect}$ 的底层标量函子。这就是标量 Whittaker 泛函能够定义谱值泛函的形式原因。
\end{note}

证明使用两个基本事实。第一，
\[\operatorname{coeff}:\DMod(\Bun_G)\longrightarrow\operatorname{Vect}\]
\srcpage{65}
经由到温和商的投影分解。第二，$\QCoh(\LocSys_{\checkG})$ 在 $\DMod(\Bun_G)^{\mathrm{temp}}$ 上的作用在所有 Satake／Hecke 算子下稳定，所以 $\operatorname{Coeff}^{\mathrm{ren}}$ 也经由温和部分分解。

\section{局部 Whittaker 范畴与温和性}\label{sec:13.2}\index{Whittaker 范畴!局部|hyperpage}\index{双 Whittaker 范畴|hyperpage}\index{温和性|hyperpage}
令 $\mathcal K=\mathbb C((t))$、$\cO=\mathbb C[[t]]$。若 DG 范畴 $\mathcal C$ 带有 $G(\mathcal K)$ 作用，则它的非退化 Whittaker 范畴为
\[\operatorname{Wh}(\mathcal C):=\mathcal C^{N(\mathcal K),\psi}.\]
在\cite{Gai18,FR23}的 Whittaker 形式体系中，局部几何 Langlands 的基本思想预言
\[\DMod(G(\mathcal K))\text{-}\mathrm{Mod}^{\mathrm{temp}}\simeq\operatorname{ShvCat}(\LocSys_{\checkG}(D^\circ)),\]
其中 $D^\circ=\Spec\mathcal K$。在这一图景下，双 Whittaker 范畴
\[\operatorname{BiWhit}(G):=\DMod((N(\mathcal K),\psi)\backslash G(\mathcal K)/(N(\mathcal K),\psi))\]
应当识别为
\[\operatorname{BiWhit}(G)\simeq\QCoh(\LocSys_{\checkG}(D^\circ)).\]
对一个 $G(\mathcal K)$-范畴 $\mathcal C$，预期的重构映射是
\[\operatorname{Wh}(\DMod(G(\mathcal K)))\otimes_{\operatorname{BiWhit}(G)}\operatorname{Wh}(\mathcal C)\longrightarrow\mathcal C.\]

\begin{transnote}{含糊点拨}
$\operatorname{BiWhit}(G)$ 是双 Whittaker 范畴，充当这里的「基环」：$\operatorname{Wh}$ 把 $G(\mathcal K)$-范畴送到它的模，$\otimes_{\operatorname{BiWhit}(G)}$ 把两边的 $\operatorname{Wh}$ 值张量起来。重构映射全忠实当且仅当 $\mathcal{C}$ 温和，这是 [Gai18] 的定理；它与 13.1 的谱定义（温和 $\leftrightarrow$ $\QCoh$）是同一件事的两种说法。[依据：原书 p.65 本段；并见 [Gai18]、[FR23]。]
\end{transnote}
重构映射全忠实，其像是由 $\operatorname{Wh}(\mathcal C)$ 生成的 $\DMod(G(\mathcal K))$-子模。因此，$\mathcal C$ 温和，当且仅当这个像是整个 $\mathcal C$。

非分歧情形更为明确。令
\[\operatorname{Sph}_{G,x}:=\DMod(G(\cO)\backslash G(\mathcal K)/G(\cO))\]
为 $x\in X$ 处的球面 Hecke 范畴。带有强 $G(\cO)$ 作用的 $G(\mathcal K)$-范畴 $\mathcal C$，受其球面不变量 $\mathcal C^{G(\cO)}$ 控制，后者是一个 $\operatorname{Sph}_{G,x}$-模。导出 Satake 给出伴随／重构图景
\[\DMod(G(\mathcal K)/G(\cO))\otimes_{\operatorname{Sph}_{G,x}}\mathcal C^{G(\cO)}\rightleftarrows\mathcal C,\]
而几何 Casselman–Shalika 等价，将从球面到 Whittaker 的比较识别为通常的从 $\operatorname{Sph}_{G,x}$ 到 $\Rep(\checkG)$ 的过渡。因此，在非分歧情形中，$\mathcal C$ 温和，当且仅当从球面不变量到 Whittaker 不变量的函子在相关的生成子范畴上保守。

对于 $\Bun_G$，每个 $x\in X$ 都通过 Hecke 修改给出 $\operatorname{Sph}_{G,x}$ 的作用。Arinkin–Gaitsgory 对 $\DMod(\Bun_G)^{\mathrm{temp}}$ 的定义，是从这一作用转到温和谱 Hecke 范畴得到的。Færgeman–Raskin 的定理证明，这一定义与 $x$ 无关\cite{FR23}。还有左温和与右温和的变体；在局部，预期它们相同，而在全局，这反映为 Whittaker 系数与所选点无关。

\section{函数论中的对应图景}\label{sec:13.3}
这一图景与经典自守理论平行。自守函数定义在
\[G(F)\backslash G(\A)/G(\cO)\]
上，而 Whittaker 系数是沿
\[(N(\A),\psi)\backslash G(\A)/G(\cO)\]
对特征标 $\psi$ 作积分。
\srcpage{66}
在几何上，后一个商被替换为一般位置定义的 $N$-约化的 Whittaker 叠，积分被替换为函子 $\operatorname{coeff}$。本章所述的保守性定理说明，限制到温和范畴后，这个几何 Whittaker 泛函能够检测所有对象的非零性。

\section{归约到幂零奇异支撑}\label{sec:13.4}\index{反温和对象|hyperpage}\index{幂零锥!非一般位置正则轨迹|hyperpage}
证明从范畴 Langlands 的受限／幂零形式所提供的比较方块开始。按此处采用的方向，示意地写为
\origdiagram{p066-comparison.pdf}{0.5829}
\begin{note}
竖直箭头应当理解为所选形式体系中可用的比较函子或投影函子，而不是同一公共范畴中的字面包含。
\end{note}
因此，只需在 $\Shv_{\mathrm{Nilp}}(\Bun_G)$ 内理解保守性的机制；在这一范畴中，可以使用微局部奇异支撑与特征圈。

回顾
\[T^*\Bun_G=\operatorname{Higgs}_G=\{(P_G,\varphi)\mid\varphi\in H^0(X,\mathfrak g_{P_G}^*\otimes\Omega_X)\}.\]
Hitchin 映射是
\[h:T^*\Bun_G\longrightarrow\cA_G:=\Gamma(X,(\mathfrak g^*/\!/G)_{\Omega_X}),\]
全局幂零锥为 $\mathrm{Nilp}:=h^{-1}(0)$。令 $\mathrm{Nilp}_{\mathrm{irreg}}\subset\mathrm{Nilp}$ 为幂零 Higgs 场一般位置正则轨迹的闭补集。因此，它的开补集就是一般位置正则幂零 Higgs 丛的轨迹。
\begin{note}
这里的“非正则”是指不在一般位置正则，并不要求在 $X$ 的每一点都非正则。
\end{note}
\begin{theorem*}{幂零层情形中的微局部定理\cite{FR25,NT25}}
作为满子范畴，有
\[\Shv_{\mathrm{Nilp}_{\mathrm{irreg}}}(\Bun_G)=\Shv_{\mathrm{Nilp}}(\Bun_G)^{\mathrm{anti}\text{-}\mathrm{temp}}.\]

\begin{transnote}{含糊点拨}
反温和指谱 Hecke 作用中负向部分（$t$-结构负次数）作用为零的层。$\mathrm{Nilp}_{\mathrm{irreg}}$ 去掉一般正则幂零轨迹，于是「支撑落在 $\mathrm{Nilp}_{\mathrm{irreg}}$」等价于「不撞该轨迹」，也就是反温和；这条微局部定理把几何条件翻译成表示论条件，后面才能把 $\operatorname{SS}(F)\subset\mathrm{Nilp}_{\mathrm{irreg}}$ 直接读成 $F$ 反温和，从而它与温和子范畴只交于零对象。[依据：原书 p.66 本段；并见 [FR25]、[NT25]。]
\end{transnote}
\end{theorem*}
它说明：若奇异支撑从不与一般位置正则幂零轨迹相交，则该对象在温和商中不可见；反之亦然。

\section{Kostant 分支与第一个系数}\label{sec:13.5}\index{Kostant 分支|hyperpage}\index{微局部指标定理|hyperpage}\index{特征圈|hyperpage}\index{微茎|hyperpage}
全局 Kostant 截面是 Chevalley 映射的 Kostant 截面 $\mathfrak g^*/\!/G\to\mathfrak g^{*,\mathrm{reg}}$ 的全局对应物，并在取 $G$ 的商后理解。进行通常的全局扭转选择后，它在 Hitchin 底空间零点上的值给出一个特殊的正则幂零轨迹，记作 $\mathfrak f^{\mathrm{glob}}\subset\mathrm{Nilp}$。令 $\mathrm{Nilp}_{\mathrm{Kos}}$ 为包含这个轨迹的唯一不可约分支。在微局部意义下，$\mathrm{Nilp}_{\mathrm{Kos}}$ 是指数 Whittaker 核的对应物，如\cite{NT25}中移位微茎的解释所示：将 Lagrange 图像 $d\psi\subset T^*\Bun_N^\Omega$ 与从 $\Bun_N^\Omega$ 到 $\Bun_G$ 的 Lagrange 对应复合，就得到全局 Kostant 切片。

特征圈计算首先对紧的可构造对象表述；可呈示范畴中的保守性陈述则由连续性与紧生成推出。

对可构造对象 $\cF\in\Shv_{\mathrm{Nilp}}(\Bun_G)$，将其特征圈写为
\[\operatorname{CC}(\cF)=\sum_{Z\subset\mathrm{Nilp}}c_Z(\cF)[Z],\qquad c_Z(\cF)\in\mathbb Z,\]
\srcpage{67}
其中求和遍及 $\mathrm{Nilp}$ 的不可约分支。
\begin{theorem*}{Færgeman–Raskin 微局部指标定理\cite{FR25}}
\[\chi(\operatorname{coeff}(\cF))=(-1)^{\dim\Bun_G}c_{\mathrm{Nilp}_{\mathrm{Kos}}}(\cF).\]
\end{theorem*}
因此，若 $\mathrm{Nilp}_{\mathrm{Kos}}$ 作为 $\operatorname{SS}(\cF)$ 的不可约分支出现（等价地，以非零特征圈重数出现），则 $\operatorname{coeff}(\cF)\ne0$。证明使用正则幂零轨迹的光滑性，以及全局 Kostant 截面沿特殊轨迹 $\mathfrak f^{\mathrm{glob}}$ 与 Hitchin 零纤维相交这一事实。

\newpage
\section{其他 Whittaker 系数与 Hecke 平移}\label{sec:13.6}\index{Hecke 移动定理|hyperpage}
令 $D=\sum_i\check\lambda_i x_i$ 为 $X$ 上取值于 $\check\Lambda^+$ 的有效除子。它对应一个对象 $V^D\in\Rep(\checkG)_{\mathrm{Ran}}$。在点互异的轨迹上，这是放置在各点 $x_i$ 的表示 $V^{\check\lambda_i}$ 的张量积；在 Ran 空间上，则是相应的因子化延拓。还有相应的 $D$-移位 Whittaker 叠；为避免固定某一种符号约定，此处将它记为 $\Bun_N^{\Omega,D}$，其系数记为 $\operatorname{coeff}_D$。
\begin{theorem*}{几何 Casselman–Shalika 公式}
\[\operatorname{coeff}_D(\cF)\simeq\operatorname{coeff}(V^D\star\cF).\]
\end{theorem*}
因此，若 $\mathrm{Nilp}_{\mathrm{Kos}}$ 作为 $\operatorname{SS}(V^D\star\cF)$ 的分支出现，则 $\operatorname{coeff}_D(\cF)\ne0$。由于 $\operatorname{Coeff}^{\mathrm{ren}}$ 与 Hecke／谱作用相容，这一非消失蕴涵 $\operatorname{Coeff}^{\mathrm{ren}}(\cF)\ne0$。

剩余的几何输入是\cite{FR25}中的 Hecke 移动定理：
\begin{theorem*}{}
若 $\operatorname{SS}(\cF)\not\subset\mathrm{Nilp}_{\mathrm{irreg}}$，则存在 $D$，使 $\mathrm{Nilp}_{\mathrm{Kos}}$ 作为 $\operatorname{SS}(V^D\star\cF)$ 的分支出现。
\end{theorem*}
该定理通过仿射 Springer 纤维的局部计算，以及一个全局铺展论证证明。非正式地说，Hecke 修改可以将任意一般位置正则的幂零奇异方向移动到 Kostant 分支。
\enlargethispage{\baselineskip}
\begin{proof}

现在设 $\cF\in\Shv_{\mathrm{Nilp}}(\Bun_G)^{\mathrm{temp}}$，并假设相应的重整化系数为零，即在受限理论中假设
\[\operatorname{Coeff}^{\mathrm{ren}}_{\mathrm{restr}}(\cF)=0.\]
Hecke 相容性于是蕴涵每个 $D$-系数都消失。由 Hecke 移动定理，$\operatorname{SS}(\cF)\subset\mathrm{Nilp}_{\mathrm{irreg}}$。因此，$\cF$ 是反温和的。由于温和子范畴与反温和子范畴的交为零，得到 $\cF=0$。这就证明了保守性。通过上面的比较方块，同一论证给出 de Rham 情形的陈述。
\end{proof}
